\subsection{Creating Conda Environments}

There are a few options for creating conda environments in OpenSARLab.
Each option come with benefits and drawbacks.


\bigskip
\centerline{\rule{13cm}{0.4pt}}
\bigskip

\subsubsection{Create Conda Environments, Register Their Kernels, and Run Any Setup Scripts in the Docker Image/s}

\paragraph{Benefits}

\begin{itemize}
\item Users don't have to create conda environments, which saves them time
\item Users don't need to know much about conda; they can just start running notebooks
\end{itemize}

\paragraph{Drawbacks}

\begin{itemize}
\item Users cannot install additional packages into their conda environments
\item Changes to environments on the docker image involve rebuilding the container and CodePipelines
\item Large environments may overrun the 20GB root volume mounted on each EC2 instance, requiring that larger, more expensive root volumes be used.
\end{itemize}


\bigskip
\centerline{\rule{13cm}{0.4pt}}
\bigskip

\subsubsection{Create Conda Environments in the Docker Image/s. Then, in the Hook Script, Sync Them to \texttt{\$Home/.local}, Register Their Kernels, and Run Any Setup Scripts}

\paragraph{Benefits}

\begin{itemize}
\item Users don't have to create conda environments, which saves them time
\item Users don't need to know much about conda at all; they can just start running notebooks
\item The environments are stored in \texttt{\$HOME/.local}, so users have permissions to install, update, remove, and debug packages\begin{itemize}
\item Environments are synced, not copied, so changes made by users will persist across server restarts
\end{itemize}
\end{itemize}

\paragraph{Drawbacks}

\begin{itemize}
\item Increases the time it takes to start an OpenSARLab server\begin{itemize}
\item Syncing environments from the docker image to \texttt{\$HOME/.local}, registering their kernels, and running any needed setup scripts all happens at server startup
\end{itemize}


\item Large environments may overrun the 20GB root volume mounted on each EC2 instance, requiring that larger, more expensive root volumes be used.
\item By storing the environment on both user volumes and EC2 node volumes, you effectively double pay for that storage.
\end{itemize}


\bigskip
\centerline{\rule{13cm}{0.4pt}}
\bigskip

\subsubsection{Leave Conda Environment Creation up to the Users}

\paragraph{Benefits}

\begin{itemize}
\item Docker images remain small, avoiding potential storage overruns on the EC2 nodes' 20GB volumes.
\item Server start ups do not require copying or syncing environments, and so require less time.
\item Users have full control over their conda environments and their changes will persist across server restarts.
\end{itemize}

\paragraph{Drawbacks}

\begin{itemize}
\item Users have to create their own conda environments\begin{itemize}
\item This requires some knowledge of conda and takes time.
\item Note: There is an \href{https://github.com/ASFOpenSARlab/opensarlab-envs}{ASF notebook repo} to aid users in building their own environments.
\end{itemize}
\end{itemize}